% !PRINTED 55
\weberpage{55}
%<<CONT
selbe, und kann der Theiler der Classe genannt werden. Sind $d, d'$ die Theiler zweier Classen $a, a'$, so ist der grösste gemeinschaftliche Theiler von $m$ und $dd'$ der Theiler der Classe $aa'$. Daraus ergiebt sich, dass die Zahlclassen, deren Zahlen relativ prim zu $m$ sind, durch Composition immer wieder solche Zahlclassen ergeben, und für diese ist dann auch die Forderung \S.\,1, 3. erfüllt.

1.\ \gesperrt{Die Zahlclassen, deren Individuen relativ prim zum Modul sind, bilden also bei der Composition durch Multiplication eine Abel'sche Gruppe.}

Diese Gruppe ist der Gegenstand unserer Betrachtungen, wobei unser Hauptziel die Bestimmung einer Basis sein soll.

Wir wollen jede Zahlclasse, die nur zu $m$ theilerfremde Zahlen enthält, mit $N$ bezeichnen, und die Gruppe der Zahlclassen $N$, deren Existenz wir jetzt nachgewiesen haben, mit $\mathfrak{N}$. Mit $n$ wollen wir jede zu $m$ theilerfremde Zahl bezeichnen.

Der Grad dieser Gruppe ist so gross, wie die Anzahl der relativen Primzahlen zu $m$, die zugleich positiv und nicht grösser als $m$ sind, und diese Zahl haben wir in \S.\,132 des ersten Bandes bestimmt. Sie ist

\weq{2}{\mu = \varphi(m) = 2^{\lambda-1} q_1^{\varkappa_1 - 1}(q_1 - 1)\, q_2^{\varkappa_2 - 1}(q_2 - 1) \ldots}

oder

\[
= q_1^{\varkappa_1 - 1}(q_1 - 1)\, q_2^{\varkappa_2 - 1}(q_2 - 1) \ldots,
\]

wenn $\lambda = 0$ ist.

Der Grad eines Elementes $N$ dieser Gruppe ist der kleinste positive Exponent $e$, zu dem man einen Repräsentanten $n$ von $N$ erheben muss, damit $n^e$ der Einheit congruent wird nach dem Modul $m$. Da jedes $e$ ein Theiler des Grades der Gruppe $\varphi(m)$ sein muss, so folgt der \gesperrt{verallgemeinerte Fermat'sche Lehrsatz}:

\weq{3}{n^{\varphi(m)} \equiv 1\ (\mathrm{mod}\ m),}

eine Formel, die für jede zu $m$ theilerfremde Zahl $n$ gilt.

Ist nun $q$ einer der ungeraden Primfactoren von $m$, und $q^\varkappa$ die höchste in $m$ aufgehende Potenz von $q$, so nehmen wir eine primitive Wurzel $g$ von $q$ an, die wir, wenn $\varkappa$ grösser als 1 ist, so wählen, dass $g^q - g$ nicht durch $q^2$ theilbar ist (Bd. I, \S.\,184). Der Kürze wegen wollen wir eine dieser letzten Bedingung genügende Zahl $g$ eine \gesperrt{primitive Wurzel von $q^2$} nennen.
